\section{Compositional Change Propagation and Version Compatibility}\label{sec:propagation}

The mechanisms developed so far treat a contract as a property of a single
method: a set of provenance- and assurance-tagged obligations against which one
implementation is synthesized and verified. Software, however, is not a flat
collection of independent methods but a graph of callers and callees, and the
durable value of a machine-checkable contract layer is realised only when that
graph is exploited. This section promotes the compositional refinement
established in Article~3 of this programme~\cite{edmonds_article3} from an
inference-time accuracy gain to a lifecycle capability. The central observation
is that because formal contracts compose along the call graph, a change to a
callee's contract mechanically \emph{recomputes} the proof obligations of every
caller, and the set of callers whose obligations can no longer be discharged is
exactly the set of callers that the change breaks---provided the compared
contracts are strong enough to express the relied-upon behaviour and the verifier
decides each obligation; where either condition fails, a caller is reported as
undetermined rather than safe, so the analysis is sound in flagging breaks but not
complete (a limitation made precise later in this section). Critically, this set is
identified by re-running the verifier, not by re-running the program: it is a
static, behavioural blast-radius analysis (\cref{fig:propagation}). The same machinery answers two
questions that the prevailing tooling treats as unrelated---``does this internal
refactor break any of my call sites?'' and ``does this dependency upgrade break
any of my call sites?''---because, at the level of contracts, an internal
refactor and a third-party version bump are the same operation. This identity is
the concrete link between the two halves of the thesis core: formal-methods
adoption and library/version compatibility are served by a single underlying
mechanism.

\begin{figure}[t]
  \centering
  \begin{tikzpicture}[
    font=\small,
    callee/.style={draw, very thick, rounded corners, align=center, fill=blue!12,
                   text width=4.0cm, minimum height=11mm, inner sep=2pt},
    caller/.style={draw, rounded corners, align=center, text width=2.0cm,
                   minimum height=9mm, inner sep=2pt},
    safe/.style={caller, fill=green!14},
    broke/.style={caller, fill=red!14},
    arr/.style={-{Latex[length=2mm]}, semithick},
  ]
    \node[broke] (A) {caller $A$\\\footnotesize \textbf{breaks}};
    \node[safe, right=9mm of A] (B) {caller $B$\\\footnotesize safe};
    \node[broke, right=9mm of B] (D) {caller $D$\\\footnotesize \textbf{breaks}};
    \node[callee, below=15mm of B] (m)
      {callee $m$\\\footnotesize contract change $(\mathit{pre}_m,\mathit{post}_m)$};

    \draw[arr] (A) -- node[left=1pt,font=\footnotesize]{calls} (m);
    \draw[arr] (B) -- (m);
    \draw[arr] (D) -- node[right=1pt,font=\footnotesize]{calls} (m);

    \node[align=center, below=4mm of m, font=\footnotesize]
      {re-verify each caller's obligations against the new contract --- \emph{without running any code}};
  \end{tikzpicture}
  \caption{Compositional change propagation as behavioural blast-radius analysis. A change to callee $m$'s contract recomputes the proof obligations of its callers; the callers whose obligations the new contract no longer discharges ($A$, $D$) are flagged statically; callers the verifier cannot decide are reported as undetermined. An internal refactor and a third-party version bump are, at the level of contracts, the same operation.}
  \label{fig:propagation}
\end{figure}

\subsection{Contracts Compose Along the Call Graph}\label{subsec:compose}

Design by contract~\cite{meyer1992dbc} casts a method specification as a
two-party agreement: the caller must establish the callee's precondition, and in
return the callee guarantees its postcondition. Modular deductive verification
operationalizes this agreement by discharging each method against the
\emph{contracts} of the methods it invokes rather than against their bodies. When
the proposed pipeline verifies a caller~$c$ that invokes a callee~$m$, OpenJML
replaces the body of~$m$ with its contract $(\mathit{pre}_m,
\mathit{post}_m)$: at the call site it asserts $\mathit{pre}_m$ as a proof
obligation on~$c$ and assumes $\mathit{post}_m$ for the continuation. The
correctness of~$c$ therefore depends on~$m$ \emph{only} through
$(\mathit{pre}_m, \mathit{post}_m)$, never through the implementation of~$m$.
This is the standard soundness story of modular verification, and it is what
makes verification scale; here it is repurposed as the engine of change
analysis. The contract is the abstraction barrier, so any change confined behind
that barrier---one that preserves $(\mathit{pre}_m, \mathit{post}_m)$---is, by
construction, invisible to every caller, and any change that crosses the
barrier---one that alters $(\mathit{pre}_m, \mathit{post}_m)$---propagates to
callers in a precisely characterizable way.

Article~3 of this programme established that the heuristic inference instrument
admits a second, weakest-precondition pass that strengthens caller contracts
using the contracts of their callees, yielding roughly $18{,}854$ additional
well-formed preconditions on the evaluation corpus beyond a single pass. That
result was reported as evidence that compositional reasoning recovers
specification content a local analysis misses. The present article reads the same
mechanism in the opposite direction. If a downstream change to
$(\mathit{pre}_m, \mathit{post}_m)$ would have altered the precondition that the
second pass computed for~$c$, then that change is, by definition, a change that
affects~$c$'s obligations. The dependency edges that the compositional pass
traverses to \emph{build} specifications are the same edges along which a change
\emph{propagates}. Recomputation is thus not new machinery but the existing
compositional pass re-evaluated under a perturbed callee contract.

\subsection{Change Propagation as Re-Verification}\label{subsec:reverify}

Let $G$ be the contract-level call graph in which an edge $c \rightarrow m$
records that caller~$c$ invokes callee~$m$ under a known set of call-site
arguments and path conditions. A proposed change replaces the callee contract
$(\mathit{pre}_m, \mathit{post}_m)$ with $(\mathit{pre}_m',
\mathit{post}_m')$. The propagation procedure is then a localized
re-verification:

\begin{enumerate}
  \item \textbf{Identify the affected frontier.} Collect the immediate callers
  of~$m$ in~$G$. Only these methods have a proof that mentions the changed
  contract directly; deeper callers are reached transitively only if a frontier
  caller's own contract changes as a result.
  \item \textbf{Recompute and re-discharge.} For each affected caller~$c$,
  re-run the modular verification of~$c$ with the callee modelled by
  $(\mathit{pre}_m', \mathit{post}_m')$. Two failure modes are
  distinguished. A \emph{precondition break} arises when $c$ established
  $\mathit{pre}_m$ at the call site but cannot establish the strengthened
  $\mathit{pre}_m'$; the verifier reports the unprovable call-site assertion. A
  \emph{postcondition break} arises when $c$'s own postcondition relied on the
  old $\mathit{post}_m$ and is no longer discharged under the weaker
  $\mathit{post}_m'$.
  \item \textbf{Propagate transitively.} If re-verification of~$c$ requires a
  change to~$c$'s own contract---for example, a strengthened precondition
  that~$c$ must now impose on \emph{its} callers---enqueue~$c$ as a changed
  callee and repeat. Because each step either discharges or strengthens a finite
  contract, and the graph is finite, the procedure terminates at a fixed point:
  the transitive set of breaking call sites together with the induced contract
  changes.
\end{enumerate}

Two properties of this procedure deserve emphasis. First, it never executes the
code. The decision of whether caller~$c$ is broken is the decision of whether an
OpenJML proof obligation is discharged by Z3, quantified over \emph{all} inputs
satisfying the path conditions, rather than over a sampled test suite. Where
test-based regression detection answers ``did any of the inputs I happened to run
change behaviour?'', contract re-verification answers ``can any input within the
specified domain now reach a violated obligation?''. Second, the analysis is
directional and asymmetric in a way that mirrors the classical theory of
behavioural subtyping and the contravariance of refinement. \emph{Weakening} a
precondition or \emph{strengthening} a postcondition is a safe change: every
caller that satisfied the old contract satisfies the new one, so no call site is
flagged, and the change is verifiably backward compatible. \emph{Strengthening}
a precondition or \emph{weakening} a postcondition is a potentially breaking
change, and the procedure reports exactly which callers it breaks and why. This
gives a developer not merely a binary compatible/incompatible verdict but a
ranked, justified work list of call sites, each accompanied by the specific
obligation that fails---the same concrete-witness signal that drives the
synthesis--verification loop of \cref{sec:verification_loop}.

\subsection{One Mechanism, Two Lifecycle Events}\label{subsec:onemech}

The procedure of \cref{subsec:reverify} is agnostic to \emph{why}
$(\mathit{pre}_m, \mathit{post}_m)$ changed. When the changed method is internal
to the codebase under development, the analysis is an impact analysis for a
refactor: a developer who tightens a guard, broadens an accepted range, or
narrows a returned set sees immediately which internal call sites must adapt.
When the changed method belongs to a third-party library and the change is
observed by comparing the contract attached to version~$v$ with that attached to
version~$v'$, the very same analysis becomes a behavioural compatibility check for
a dependency upgrade. The question ``which of my $N$ call sites does this upgrade
break behaviourally?'' is thereby promoted to a first-class, machine-answerable
lifecycle event.

This framing addresses a well-documented and costly gap in current dependency
practice. Semantic versioning~\cite{preston2013semver} is the dominant
convention for signalling compatibility, yet large-scale studies of Maven
Central repeatedly find that version numbers are unreliable predictors of actual
breakage: a substantial fraction of releases labelled non-breaking introduce
incompatible changes~\cite{raemaekers2017semver,ochoa2022breaking}, and
dependency networks evolve fast enough that clients are exposed to such changes
continually~\cite{decan2019empirical}. The tooling that does exist operates
predominantly at the level of \emph{signatures}---method removal, parameter-type
changes, visibility reductions---as exemplified by APIDiff and related
detectors~\cite{brito2018apidiff,jezek2017magic}. Signature-level analysis,
however, is blind to \emph{behavioural} breaking changes: a method whose
signature is untouched but whose contract has shifted, for instance one that
begins rejecting an input it previously accepted or that stops guaranteeing a
postcondition a caller relied upon. Such behavioural incompatibilities are both
prevalent and damaging to clients~\cite{mostafa2017behavioral,jayasuriya2024understanding},
and static checking of library updates has been pursued precisely because
recompilation alone does not surface them~\cite{foo2018efficient}. Recent work
even applies LLMs to \emph{repair} clients broken by dependency
updates~\cite{islam2025byam}; the contract-propagation analysis proposed here is
complementary, supplying the precise, per-call-site diagnosis of \emph{what}
broke and \emph{why} that such a repair step requires as input. By making the
unit of comparison the formal contract rather than the signature, the pipeline
detects exactly the behavioural incompatibilities that semantic-versioning labels
and signature diffs miss, and it does so soundly over the input domain rather
than over a regression sample. The contract attached to a compiled
artefact---the embedding mechanism validated in Article~2 of this
programme~\cite{edmonds_article2}---is what makes the library side of this
comparison feasible without library source, and the diff-only, cache-backed CI
integration of Article~4~\cite{edmonds_article4} is what makes re-running the
analysis on every upgrade economical.

\subsection{A Worked Version Step}\label{subsec:worked}

Consider a utility library that, at version~$v$, exposes an integer-parsing
helper modelled by the contract in \cref{lst:callee-old}. The method accepts any
string and returns a sentinel of $-1$ when the argument is not a valid integer;
its postcondition records both the success and the sentinel cases.

\begin{lstlisting}[caption={Library method \texttt{toInt} at version~$v$ (callee contract).},label={lst:callee-old}]
//@ requires s != null;
//@ ensures isParsable(s)  ==> \result == parseInt(s);
//@ ensures !isParsable(s) ==> \result == -1;
public static int toInt(String s) { ... }
\end{lstlisting}

A downstream client contains the call site in \cref{lst:caller}. Its own
postcondition---that the returned count is non-negative---is discharged against
version~$v$ because every branch of \texttt{toInt} returns either a parsed value
guarded by the caller or the sentinel~$-1$, which the caller maps to~$0$.

\begin{lstlisting}[caption={Client call site relying on the version-$v$ contract.},label={lst:caller}]
//@ ensures \result >= 0;
public int countOrZero(String s) {
    int n = LibUtil.toInt(s);   // callee contract assumed here
    return n < 0 ? 0 : n;
}
\end{lstlisting}

At version~$v'$ the library maintainer revises \texttt{toInt} to throw on
malformed input rather than return a sentinel---a change the maintainer regards
as a clean-up and ships under a non-breaking version label. The revised contract
is shown in \cref{lst:callee-new}: the precondition is \emph{strengthened} to
demand a parsable argument, and the success postcondition is retained while the
sentinel clause is removed.

\begin{lstlisting}[caption={Library method \texttt{toInt} at version~$v'$ (strengthened precondition).},label={lst:callee-new}]
//@ requires s != null && isParsable(s);
//@ ensures \result == parseInt(s);
public static int toInt(String s) { ... }
\end{lstlisting}

Running the procedure of \cref{subsec:reverify} with
$(\mathit{pre}_{\mathtt{toInt}}, \mathit{post}_{\mathtt{toInt}})$ replaced by
the version-$v'$ contract recomputes the obligations of \texttt{countOrZero}.
The strengthened precondition $\mathtt{isParsable(s)}$ becomes a call-site proof
obligation that \texttt{countOrZero} cannot discharge: nothing in the caller
establishes that \texttt{s} is parsable, so the verifier reports a precondition
break at the call site. The analysis therefore flags \texttt{countOrZero} as a
behaviourally broken caller of the upgrade, names the unprovable obligation
$\mathtt{isParsable(s)}$ as the cause, and---because no contract change to
\texttt{countOrZero} would discharge the obligation without either adding a guard
or strengthening \texttt{countOrZero}'s own precondition---surfaces the two
candidate remediations as the induced contract change to be propagated to
\emph{its} callers. A signature-level tool would report no break whatsoever,
since the signature of \texttt{toInt} is unchanged; the contract-level analysis
correctly identifies the behavioural incompatibility that the non-breaking version
label conceals. \Cref{tab:propagation-outcomes} summarises how the four
contract-change directions classify under the analysis.

\begin{table}[t]
  \centering
  \caption{Classification of callee-contract changes by the re-verification
  procedure. Safe changes flag no callers; breaking changes yield a justified
  per-call-site work list.}
  \label{tab:propagation-outcomes}
  \begin{tabular}{@{}lll@{}}
    \toprule
    Change to callee contract & Compatibility & Caller effect \\
    \midrule
    Weaken precondition    & Safe     & No call site flagged \\
    Strengthen postcondition & Safe   & No call site flagged \\
    Strengthen precondition & Breaking & Call-site assertion unprovable \\
    Weaken postcondition    & Breaking & Caller postcondition undischarged \\
    \bottomrule
  \end{tabular}
\end{table}

\subsection{Scope, Frontier, and Honest Limits}\label{subsec:proplimits}

The propagation analysis inherits both the strengths and the limits of the
underlying instrument. Its soundness is the soundness of OpenJML extended static
checking discharged by Z3: a caller reported as compatible is compatible over the
entire specified input domain, not merely over a tested sample, but a caller for
which Z3 times out---on the multiplicative reasoning, array theory, or
rich-precondition cases known to defeat the solver---must be reported under the
tiered-assurance scheme as \emph{undetermined} rather than silently passed. The
analysis is also only as good as the contracts it compares. For first-party code
the contracts are those the pipeline already maintains; for third-party libraries
the contract attached to each version is, in the brownfield setting, an
\emph{inferred} contract recovered from the library's code, and an inferred
callee contract skews weak, so a genuine behavioural break may be missed when the
inferred contract failed to capture the guarantee a caller actually relied upon.
This is the same weak-specification limitation that motivates ratification and
inherited defaults elsewhere in the pipeline, and it confines the trustworthy
application of version-compatibility analysis to the API surface and the modules
within the specified frontier rather than to an entire legacy dependency tree.
Within that frontier, however, the analysis delivers what semantic-versioning
labels and signature diffs cannot: a sound, behaviour-level, per-call-site verdict
on whether a refactor or an upgrade is safe, computed without running a single
test. Establishing the empirical precision and recall of this analysis against a
ground-truth corpus of known behavioural breaking changes is left to the
evaluation; the architectural claim here is that the capability falls directly
out of compositional contracts and requires no machinery beyond the
re-verification of an already-built specification graph.
